% GENERATED FILE. Edit canonical lessons, then run npm run generate:books.
% canonical-source-hash: fnv1a64:67e217315b956359
% canonical-lessons: ES-C46-le

\chapter{Le --- To Him, To Her, To You}
\label{ch:pronoun-le}
\hlchaptermodality{93}

\begin{chapteropening}
\textbf{By the end of this chapter:} \emph{I can say le digo, and know that le does not mark gender.}

Complete ES-C46-le: produce le digo, say how many genders le has, and name the Latin case it comes from.
\end{chapteropening}

\section[le]{le — \textquotedblleft{}to him,\textquotedblright{} \textquotedblleft{}to her,\textquotedblright{} \textquotedblleft{}to you\textquotedblright{}}

\label{lesson:ES-C46-le}



\begin{quote}
\emph{Lo} and \emph{la} answer the question \textbf{what}. This chapter answers a different one: \textbf{to whom}.
\end{quote}

\begin{grammarlens}[title={Grammar Lens: the other question}]
\begin{quote}
\textbf{le} — \emph{to him}, \emph{to her}, \emph{to you} (with \emph{usted}).
\end{quote}

\begin{quote}
\emph{\textbf{Le} digo.} — I tell \textbf{him}. Or her. Or you.
\end{quote}

Some verbs land on a thing: you have it, you make it, you want it. Others land on a \textbf{person you are aiming at} — you tell them, you speak to them, you give them something. Spanish uses a different pronoun for the second kind, and \emph{le} is it.

Same slot. Still in front of the verb. Nothing about the shape of a sentence has changed in eleven chapters.

\textbf{And here is the relief:} \emph{le} does \textbf{not} mark gender.

\begin{quote}
\emph{Le digo.} — I tell him. \textbf{and} I tell her. \textbf{and} I tell you.
\end{quote}

After four chapters of choosing between \emph{lo} and \emph{la} by the gender of a noun, this one asks you for nothing. One form covers three people. If you need to say which, you add it in words — but the pronoun itself does not care.
\end{grammarlens}

\begin{cousinweb}
\emph{Lo} and \emph{le} are the same Latin word again — and this time the difference is not wear, it is \textbf{grammar}.

Latin changed a word's ending to show its job in the sentence. \emph{Ille}, \textquotedblleft{}that one,\textquotedblright{} had one form for the thing an action hit and another for the person it was aimed at:

\begin{itemize}
\raggedright
  \item \textbf{illum} — the one the verb acts on $\to$ Spanish \textbf{lo}.
  \item \textbf{illī} — the one it is aimed at $\to$ Spanish \textbf{le}.
\end{itemize}

Latin called these the accusative and the dative. Almost all of that system collapsed on the way to Spanish; nouns lost their endings entirely. But in these small, constantly-used pronouns, \textbf{two of the cases survived}, and you are looking at them.

English had the same distinction and lost it more completely. It survives in a word order rather than a shape: \emph{I told \textbf{him} the story} — \emph{him} is doing exactly what \emph{le} does, and English simply stopped marking it.
\end{cousinweb}

\subsection*{Your turn}
Say these aloud:

\begin{itemize}
\raggedright
  \item \textquotedblleft{}Lo tengo\textquotedblright{} then \textquotedblleft{}Le digo\textquotedblright{} — what, then to whom
  \item \textquotedblleft{}Le digo\textquotedblright{} three times, meaning him, then her, then usted
\end{itemize}

\begin{tcolorbox}[breakable,colback=teal!4,colframe=teal!35!black,title={Before you move on}]
Say \textquotedblleft{}I tell him.\textquotedblright{} (\emph{\textbf{Le digo}}.) And \textquotedblleft{}I tell her.\textquotedblright{} (\emph{\textbf{Le digo}} — the same.) Which Latin case is \emph{le}? (The \textbf{dative}, \emph{illī}.) Next: its plural, and a discovery about how little of this system is actually new.
\end{tcolorbox}
